%!TeX root=../emmatop.tex
\chapter[Chapter \thechapter]{}
\lettrine[lines=4,lraise=0.3]{S}{mall} heart had Harriet for visiting. Only half an hour before her friend called for her at Mrs Goddard's, her evil stars had led her to the very spot where, at that moment, a trunk, directed to The Rev. Philip Elton, White-Hart, Bath, was to be seen under the operation of being lifted into the butcher's cart, which was to convey it to where the coaches past; and every thing in this world, excepting that trunk and the direction, was consequently a blank.

She went, however; and when they reached the farm, and she was to be put down, at the end of the broad, neat gravel walk, which led between espalier apple-trees to the front door, the sight of every thing which had given her so much pleasure the autumn before, was beginning to revive a little local agitation; and when they parted, Emma observed her to be looking around with a sort of fearful curiosity, which determined her not to allow the visit to exceed the proposed quarter of an hour. She went on herself, to give that portion of time to an old servant who was married, and settled in Donwell.

The quarter of an hour brought her punctually to the white gate again; and Miss Smith receiving her summons, was with her without delay, and unattended by any alarming young man. She came solitarily down the gravel walk—a Miss Martin just appearing at the door, and parting with her seemingly with ceremonious civility.

\begin{figure}[tbph]
\centering
\includegraphics[width=.9\linewidth]{23room}
\caption{In that very room she had been measured}
\end{figure}

Harriet could not very soon give an intelligible account. She was feeling too much; but at last Emma collected from her enough to understand the sort of meeting, and the sort of pain it was creating. She had seen only Mrs Martin and the two girls. They had received her doubtingly, if not coolly; and nothing beyond the merest commonplace had been talked almost all the time—till just at last, when Mrs Martin's saying, all of a sudden, that she thought Miss Smith was grown, had brought on a more interesting subject, and a warmer manner. In that very room she had been measured last September, with her two friends. There were the pencilled marks and memorandums on the wainscot by the window. He had done it. They all seemed to remember the day, the hour, the party, the occasion—to feel the same consciousness, the same regrets—to be ready to return to the same good understanding; and they were just growing again like themselves, (Harriet, as Emma must suspect, as ready as the best of them to be cordial and happy,) when the carriage reappeared, and all was over. The style of the visit, and the shortness of it, were then felt to be decisive. Fourteen minutes to be given to those with whom she had thankfully passed six weeks not six months ago!—Emma could not but picture it all, and feel how justly they might resent, how naturally Harriet must suffer. It was a bad business. She would have given a great deal, or endured a great deal, to have had the Martins in a higher rank of life. They were so deserving, that a little higher should have been enough: but as it was, how could she have done otherwise?—Impossible!—She could not repent. They must be separated; but there was a great deal of pain in the process—so much to herself at this time, that she soon felt the necessity of a little consolation, and resolved on going home by way of Randalls to procure it. Her mind was quite sick of Mr Elton and the Martins. The refreshment of Randalls was absolutely necessary.

It was a good scheme; but on driving to the door they heard that neither »master nor mistress was at home;« they had both been out some time; the man believed they were gone to Hartfield.

»This is too bad,« cried Emma, as they turned away. »And now we shall just miss them; too provoking!—I do not know when I have been so disappointed.« And she leaned back in the corner, to indulge her murmurs, or to reason them away; probably a little of both—such being the commonest process of a not ill-disposed mind. Presently the carriage stopt; she looked up; it was stopt by Mr and Mrs Weston, who were standing to speak to her. There was instant pleasure in the sight of them, and still greater pleasure was conveyed in sound—for Mr Weston immediately accosted her with,

»How d'ye do?—how d'ye do?—We have been sitting with your father—glad to see him so well. Frank comes to-morrow—I had a letter this morning—we see him to-morrow by dinner-time to a certainty—he is at Oxford to-day, and he comes for a whole fortnight; I knew it would be so. If he had come at Christmas he could not have staid three days; I was always glad he did not come at Christmas; now we are going to have just the right weather for him, fine, dry, settled weather. We shall enjoy him completely; every thing has turned out exactly as we could wish.«

There was no resisting such news, no possibility of avoiding the influence of such a happy face as Mr Weston's, confirmed as it all was by the words and the countenance of his wife, fewer and quieter, but not less to the purpose. To know that she thought his coming certain was enough to make Emma consider it so, and sincerely did she rejoice in their joy. It was a most delightful reanimation of exhausted spirits. The worn-out past was sunk in the freshness of what was coming; and in the rapidity of half a moment's thought, she hoped Mr Elton would now be talked of no more.

Mr Weston gave her the history of the engagements at Enscombe, which allowed his son to answer for having an entire fortnight at his command, as well as the route and the method of his journey; and she listened, and smiled, and congratulated.

»I shall soon bring him over to Hartfield,« said he, at the conclusion.

Emma could imagine she saw a touch of the arm at this speech, from his wife.

»We had better move on, Mr Weston,« said she, »we are detaining the girls.«

»Well, well, I am ready;«—and turning again to Emma, »but you must not be expecting such a very fine young man; you have only had my account you know; I dare say he is really nothing extraordinary:«—though his own sparkling eyes at the moment were speaking a very different conviction.

Emma could look perfectly unconscious and innocent, and answer in a manner that appropriated nothing.

»Think of me to-morrow, my dear Emma, about four o'clock,« was Mrs Weston's parting injunction; spoken with some anxiety, and meant only for her.

»Four o'clock!—depend upon it he will be here by three,« was Mr Weston's quick amendment; and so ended a most satisfactory meeting. Emma's spirits were mounted quite up to happiness; every thing wore a different air; James and his horses seemed not half so sluggish as before. When she looked at the hedges, she thought the elder at least must soon be coming out; and when she turned round to Harriet, she saw something like a look of spring, a tender smile even there.

»Will Mr Frank Churchill pass through Bath as well as Oxford?«—was a question, however, which did not augur much.

But neither geography nor tranquillity could come all at once, and Emma was now in a humour to resolve that they should both come in time.

The morning of the interesting day arrived, and Mrs Weston's faithful pupil did not forget either at ten, or eleven, or twelve o'clock, that she was to think of her at four.

»My dear, dear anxious friend,«—said she, in mental soliloquy, while walking downstairs from her own room, »always overcareful for every body's comfort but your own; I see you now in all your little fidgets, going again and again into his room, to be sure that all is right.« The clock struck twelve as she passed through the hall. »'Tis twelve; I shall not forget to think of you four hours hence; and by this time to-morrow, perhaps, or a little later, I may be thinking of the possibility of their all calling here. I am sure they will bring him soon.«

She opened the parlour door, and saw two gentlemen sitting with her father—Mr Weston and his son. They had been arrived only a few minutes, and Mr Weston had scarcely finished his explanation of Frank's being a day before his time, and her father was yet in the midst of his very civil welcome and congratulations, when she appeared, to have her share of surprize, introduction, and pleasure.

The Frank Churchill so long talked of, so high in interest, was actually before her—he was presented to her, and she did not think too much had been said in his praise; he was a very good looking young man; height, air, address, all were unexceptionable, and his countenance had a great deal of the spirit and liveliness of his father's; he looked quick and sensible. She felt immediately that she should like him; and there was a well-bred ease of manner, and a readiness to talk, which convinced her that he came intending to be acquainted with her, and that acquainted they soon must be.

He had reached Randalls the evening before. She was pleased with the eagerness to arrive which had made him alter his plan, and travel earlier, later, and quicker, that he might gain half a day.

»I told you yesterday,« cried Mr Weston with exultation, »I told you all that he would be here before the time named. I remembered what I used to do myself. One cannot creep upon a journey; one cannot help getting on faster than one has planned; and the pleasure of coming in upon one's friends before the look-out begins, is worth a great deal more than any little exertion it needs.«

»It is a great pleasure where one can indulge in it,« said the young man, »though there are not many houses that I should presume on so far; but in coming home I felt I might do any thing.«

The word home made his father look on him with fresh complacency. Emma was directly sure that he knew how to make himself agreeable; the conviction was strengthened by what followed. He was very much pleased with Randalls, thought it a most admirably arranged house, would hardly allow it even to be very small, admired the situation, the walk to Highbury, Highbury itself, Hartfield still more, and professed himself to have always felt the sort of interest in the country which none but one's own country gives, and the greatest curiosity to visit it. That he should never have been able to indulge so amiable a feeling before, passed suspiciously through Emma's brain; but still, if it were a falsehood, it was a pleasant one, and pleasantly handled. His manner had no air of study or exaggeration. He did really look and speak as if in a state of no common enjoyment.

Their subjects in general were such as belong to an opening acquaintance. On his side were the inquiries,—»Was she a horsewoman?—Pleasant rides?—Pleasant walks?—Had they a large neighbourhood?—Highbury, perhaps, afforded society enough?—There were several very pretty houses in and about it.—Balls—had they balls?—Was it a musical society?«

But when satisfied on all these points, and their acquaintance proportionably advanced, he contrived to find an opportunity, while their two fathers were engaged with each other, of introducing his mother-in-law, and speaking of her with so much handsome praise, so much warm admiration, so much gratitude for the happiness she secured to his father, and her very kind reception of himself, as was an additional proof of his knowing how to please—and of his certainly thinking it worth while to try to please her. He did not advance a word of praise beyond what she knew to be thoroughly deserved by Mrs Weston; but, undoubtedly he could know very little of the matter. He understood what would be welcome; he could be sure of little else. »His father's marriage,« he said, »had been the wisest measure, every friend must rejoice in it; and the family from whom he had received such a blessing must be ever considered as having conferred the highest obligation on him.«

He got as near as he could to thanking her for Miss Taylor's merits, without seeming quite to forget that in the common course of things it was to be rather supposed that Miss Taylor had formed Miss Woodhouse's character, than Miss Woodhouse Miss Taylor's. And at last, as if resolved to qualify his opinion completely for travelling round to its object, he wound it all up with astonishment at the youth and beauty of her person.

»Elegant, agreeable manners, I was prepared for,« said he; »but I confess that, considering every thing, I had not expected more than a very tolerably well-looking woman of a certain age; I did not know that I was to find a pretty young woman in Mrs Weston.«

»You cannot see too much perfection in Mrs Weston for my feelings,« said Emma; »were you to guess her to be eighteen, I should listen with pleasure; but she would be ready to quarrel with you for using such words. Don't let her imagine that you have spoken of her as a pretty young woman.«

»I hope I should know better,« he replied; »no, depend upon it, (with a gallant bow,) that in addressing Mrs Weston I should understand whom I might praise without any danger of being thought extravagant in my terms.«

Emma wondered whether the same suspicion of what might be expected from their knowing each other, which had taken strong possession of her mind, had ever crossed his; and whether his compliments were to be considered as marks of acquiescence, or proofs of defiance. She must see more of him to understand his ways; at present she only felt they were agreeable.

She had no doubt of what Mr Weston was often thinking about. His quick eye she detected again and again glancing towards them with a happy expression; and even, when he might have determined not to look, she was confident that he was often listening.

Her own father's perfect exemption from any thought of the kind, the entire deficiency in him of all such sort of penetration or suspicion, was a most comfortable circumstance. Happily he was not farther from approving matrimony than from foreseeing it.—Though always objecting to every marriage that was arranged, he never suffered beforehand from the apprehension of any; it seemed as if he could not think so ill of any two persons' understanding as to suppose they meant to marry till it were proved against them. She blessed the favouring blindness. He could now, without the drawback of a single unpleasant surmise, without a glance forward at any possible treachery in his guest, give way to all his natural kind-hearted civility in solicitous inquiries after Mr Frank Churchill's accommodation on his journey, through the sad evils of sleeping two nights on the road, and express very genuine unmixed anxiety to know that he had certainly escaped catching cold—which, however, he could not allow him to feel quite assured of himself till after another night.

A reasonable visit paid, Mr Weston began to move.—»He must be going. He had business at the Crown about his hay, and a great many errands for Mrs Weston at Ford's, but he need not hurry any body else.« His son, too well bred to hear the hint, rose immediately also, saying,

»As you are going farther on business, sir, I will take the opportunity of paying a visit, which must be paid some day or other, and therefore may as well be paid now. I have the honour of being acquainted with a neighbour of yours, (turning to Emma,) a lady residing in or near Highbury; a family of the name of Fairfax. I shall have no difficulty, I suppose, in finding the house; though Fairfax, I believe, is not the proper name—I should rather say Barnes, or Bates. Do you know any family of that name?«

»To be sure we do,« cried his father; »Mrs Bates—we passed her house—I saw Miss Bates at the window. True, true, you are acquainted with Miss Fairfax; I remember you knew her at Weymouth, and a fine girl she is. Call upon her, by all means.«

»There is no necessity for my calling this morning,« said the young man; »another day would do as well; but there was that degree of acquaintance at Weymouth which\longdash«

»Oh! go to-day, go to-day. Do not defer it. What is right to be done cannot be done too soon. And, besides, I must give you a hint, Frank; any want of attention to her here should be carefully avoided. You saw her with the Campbells, when she was the equal of every body she mixed with, but here she is with a poor old grandmother, who has barely enough to live on. If you do not call early it will be a slight.«

The son looked convinced.

»I have heard her speak of the acquaintance,« said Emma; »she is a very elegant young woman.«

He agreed to it, but with so quiet a »Yes,« as inclined her almost to doubt his real concurrence; and yet there must be a very distinct sort of elegance for the fashionable world, if Jane Fairfax could be thought only ordinarily gifted with it.

»If you were never particularly struck by her manners before,« said she, »I think you will to-day. You will see her to advantage; see her and hear her—no, I am afraid you will not hear her at all, for she has an aunt who never holds her tongue.«

»You are acquainted with Miss Jane Fairfax, sir, are you?« said Mr Woodhouse, always the last to make his way in conversation; »then give me leave to assure you that you will find her a very agreeable young lady. She is staying here on a visit to her grandmama and aunt, very worthy people; I have known them all my life. They will be extremely glad to see you, I am sure; and one of my servants shall go with you to shew you the way.«

»My dear sir, upon no account in the world; my father can direct me.«

»But your father is not going so far; he is only going to the Crown, quite on the other side of the street, and there are a great many houses; you might be very much at a loss, and it is a very dirty walk, unless you keep on the footpath; but my coachman can tell you where you had best cross the street.«

Mr Frank Churchill still declined it, looking as serious as he could, and his father gave his hearty support by calling out, »My good friend, this is quite unnecessary; Frank knows a puddle of water when he sees it, and as to Mrs Bates's, he may get there from the Crown in a hop, step, and jump.«

They were permitted to go alone; and with a cordial nod from one, and a graceful bow from the other, the two gentlemen took leave. Emma remained very well pleased with this beginning of the acquaintance, and could now engage to think of them all at Randalls any hour of the day, with full confidence in their comfort.